\section{Introduction}
\label{sec:introduction}

Consider an agent entering a previously unseen Web application. Rather than
receiving a single user goal, the agent must determine what the application
can do and organize its experience into a functional model that can support
future tasks. Such a model could amortize interaction across tasks: instead of
rediscovering the interface for every new request, a downstream agent could
reuse grounded knowledge about available functions, where they apply, and what
observable outcomes they produce. Building this model, however, requires more
than passively reading a page. Many application functions become evident only
after the agent changes state through interaction.

This setting motivates \emph{open-ended Web exploration}
\citep{xie2025guiexplorer,nica2025uiexplore,fan2025guibee}. Task-conditioned
agents typically search for a path that completes a specified objective; they
need not discover functionality irrelevant to that objective or preserve it
for later use. An open-ended explorer instead proposes its own functional
hypotheses, visits different application states, and accumulates knowledge for
goals that are not yet known. These paradigms are complementary rather than
universally ordered: open-ended exploration is aligned with reusable model
construction, but its autonomy removes constraints supplied by a concrete goal.

The first resulting challenge is epistemic. A proposed function is only a
candidate, and an executor-reported success establishes only that an
interaction was completed. Neither establishes that the intended
application-level outcome occurred. Persisting either signal as fact can
introduce unsupported knowledge that later misleads planning or verification.
Functional model induction therefore needs an explicit admission boundary:
the agent should state an expected observable outcome before acting, preserve
the resulting interaction evidence, and admit the claim only when that
evidence supports the frozen expectation.

The second challenge concerns the process used to obtain that evidence
\citep{zheng2025webguard,sun2026ossentinel,mohammadmirzaei2026osguard}.
Testing an unknown function may require clicking, entering data, submitting a
form, or confirming a transition. Depending on the visible application state,
such an action may delete information, send content, change access, enter a
transactional workflow, or prepare another consequential operation. Risk is
therefore not a context-free property of a function name: it depends on the
current observation and the particular browser action about to be executed.
Moreover, a function can be learned correctly even though the interaction used
to learn it was consequential. Trustworthy model induction must consequently
make both the evidential status of acquired knowledge and the environmental
risk of acquiring it explicit.

We study open-ended exploration as \emph{evidence production under dual risk}
and introduce \textbf{VERA}, a framework for Verification and Environmental
Risk Awareness in functional model induction. VERA proposes a
location-conditioned functional claim together with a frozen expected
observable outcome. Before each selected browser action, it records a
context-conditioned risk judgment grounded in the current GUI observation.
After execution, it links the action, executor status, and before--after
observations to the claim, and admits the claim to downstream-usable knowledge
only when the evidence supports the expected outcome. Persistent unresolved
hypotheses and replay recover further opportunities for evidence production.
The contribution is the protocol connecting these records into a traceable
functional model.

Controlled studies support these design choices. Evidence-grounded admission
attains 96.97\% precision while retaining 96.97\% of human-supported functional
knowledge. Under the same ordinary-attempt budget, VERA's full exploration
mechanism improves website-macro-average evidence-supported functional coverage
by 11.9 percentage points over linear exploration, with replay costs reported
separately. Finally, adding visual context improves acceptable risk-type
accuracy by 35.4 percentage points on an external context-challenge set, while
binary detection exhibits a higher-recall, lower-precision trade-off. Together,
these results provide bounded evidence for reliable knowledge admission,
sustained evidence production, and context-grounded risk awareness. The
evaluated risk records do not alter the selected action.

Our contributions are:
\begin{itemize}
    \item We formulate open-ended functional model induction as evidence
    production under epistemic and environmental interaction risks, and define
    a traceable model that separates candidate discovery, executor-reported
    success, evidence judgment, and knowledge admission.
    \item We show how outcome verification can govern persistent knowledge
    admission, while a persistent frontier and replay sustain evidence
    production for unresolved claims.
    \item We make evidence acquisition traceable by linking each selected
    action's context-conditioned risk record to its claim, execution trace, and
    outcome evidence, and evaluate this record on workflow and external cases.
\end{itemize}
